\section{Successive faces and projective windows}\label{win:section}

Throughout this section let $Y$ be a connected smooth projective variety
and let $V\subset Y$ be a nonempty connected semialgebraic open subset
that is locally pseudoconvex at its boundary.  A domain is connected.
A matrix inequality means
positive definiteness of the indicated real symmetric matrix.  Empty matrices
are permitted; thus the test associated with an empty lower block is simply
$\Im s>0$.  The open kernel always has the neighborhood meaning of
Definition~\ref{prep:kernel}.  In particular, its complement is the limiting
accumulation set of the forbidden points, and not merely the set of pointwise
failures of eventual membership.

We use the following precise consequences of the preceding preparatory
results.  Specializations with fixed rational weights preserve local
pseudoconvexity and have semialgebraic open kernels; after restriction at
generic real centers their tests are invariant under real translations.
The references are Lemmas~\ref{prep:kernel-pseudoconvexity}
and~\ref{prep:finite-jet-specialization} and
Corollary~\ref{prep:generic-freezing}.
We also use the projective hole tests with a container defined on a full
parameter box, including its real edge: boundary trapping, interior trapping
on a nonpluripolar slice, and propagation through a connected projection
component, as in Lemmas~\ref{holes:boundary-trapping}
and~\ref{holes:interior-trapping} and
Corollary~\ref{holes:container-propagation}.
Whenever these results are applied below, the projective family is smooth
and proper, its fibres are connected, its evaluation is holomorphic, and its
open subset is the inverse image of the locally pseudoconvex domain under
consideration.  No conclusion about the boundary behavior of arbitrary
holomorphic functions is part of these inputs.

\subsection{Algebraic plaques and the contact quotient}

\begin{lemma}[Algebraicity of the null plaques]\label{win:plaques}
Let $M$ be a regular patch of a real algebraic hypersurface in a smooth
complex algebraic $n$-fold.  Suppose one side is locally pseudoconvex and
the Levi form has constant rank $r$ on $M$.  Put $\ell=n-r-1$.
The Levi-null distribution integrates into complex $\ell$-dimensional
plaques.  Each sufficiently small plaque is open in an irreducible complex
algebraic variety of dimension $\ell$.  Locally at generic transverse
parameters these varieties form a real analytic family of bounded degree.
After complexifying the parameters, taking the relevant components, and
resolving, they have smooth projective models in a smooth proper family.
\end{lemma}

\begin{proof}
Choose a real analytic defining function $\rho$, and let $\theta$ be its
characteristic one-form on $M$.  Write $D=\ker\theta=TM\cap JTM$.
Since the Levi form is semidefinite, its null space agrees with the
characteristic distribution
\[
 K=\{v\in D:d\theta(v,D)=0\}.
\]
This distribution has real rank $2\ell$ and is $J$-invariant.  If $X$ is a
section of $K$ and $Y$ is a section of $D$, then
$\theta([X,Y])=-d\theta(X,Y)=0$.  Applying the Jacobi identity to two
sections of $K$ and one section of $D$ shows that their bracket again has
this property.  Consequently $K$ is involutive.  The analytic Frobenius
theorem gives analytic leaves, and their $J$-invariant tangent spaces make
them complex submanifolds.

In affine algebraic coordinates write the real polynomial defining $M$ as
\[
             \rho(z,\bar z)=v(z)^*C v(z),\qquad C=C^*,
\]
where $v$ is a finite vector of holomorphic monomials; adjoining a constant
monomial is harmless.  For a parametrization $f$ of a complex plaque,
polarization of the identity $\rho(f(\zeta),\overline{f(\zeta)})=0$ gives
\[
             v(f(\eta))^*C v(f(\zeta))=0
\]
for all $\eta,\zeta$ in a smaller plaque.  Indeed the polarized expression
is holomorphic in $\zeta$ and antiholomorphic in $\eta$, and its restriction
to the diagonal determines every coefficient of its convergent series.
The complex span $W$ of the vectors $v(f(\zeta))$ is therefore totally
isotropic for $C$.

The algebraic set $A=v^{-1}(W)$ contains the plaque and lies in
$\{\rho=0\}$.  Every complex submanifold of $M$ has tangent space in $K$:
the restriction of $\rho$ vanishes identically, so its Levi form vanishes
there, and semidefiniteness identifies its zero vectors with its kernel.
Applying this at regular points of each local component of $A$ bounds
that component's dimension by $\ell$.  The component containing the
plaque thus has dimension exactly $\ell$, and the plaque is open in it.

Only finitely many evaluation vectors are needed to span $W$.  Choose
them by fixed evaluations in analytic plaque coordinates, and restrict
to a parameter neighborhood where their rank is maximal and constant.
Their span is then an analytic map into a Grassmannian.  The equations
$v(z)\in W$ have uniformly bounded degrees, and their projective closures
therefore have bounded degree.  Track the component containing the marked
plaque, using its irreducible germ.  The family and its evaluation have
the analytic algebraic presentation of Lemma~\ref{prep:projective-models}.
Complexification, a local branch in the component space,
relative resolution, and generic smoothness give the stated smooth
projective family.  The generic restrictions here remove rank-drop and
component-change loci; they do not assert uniformity over their closures.
\end{proof}

\begin{lemma}[The period matrix]\label{win:period}
In the situation of Lemma~\ref{win:plaques}, there is a real analytic
$(r+1)$-parameter family of plaques, with coordinates $(s,w)\in
\R\times\R^r$, whose complexification evaluates with full complex rank
near the original plaque.  The remaining $r$ real transverse coordinates
can be retained as external parameters.  At a fixed real center $x_1$,
let $L$ be the smooth projective model of its plaque.  There is a symmetric
meromorphic $r\times r$ matrix $H$ on $L$, holomorphic near the original
plaque, such that $\Im H>0$ there and the weighted tangent inequality is
\begin{equation}\label{win:quadratic-test}
       q_z(s,w)=\Im s-(\Im w)^{\mathsf T}
                          (\Im H(z))^{-1}\Im w>0.
\end{equation}
The meromorphic data depend analytically on the generic external
parameters.
\end{lemma}

\begin{proof}
The quotient of $M$ by the characteristic plaques is a real analytic
contact manifold of dimension $2r+1$.  Indeed the hyperplane field $D$
descends because $[K,D]\subset D$, and its induced exterior derivative is
nondegenerate.  Choose Darboux coordinates with contact form
$da-\sum_j v_j\,du_j$.  Holding $v$ fixed and writing
$a=s+v^{\mathsf T}w$, $u=w$, gives a family on which the pulled-back
contact form is $ds$.  Thus the $w$-directions form a Legendrian
Lagrangian and the $s$-direction is transverse.  The parameters $v$ are
the external parameters just mentioned.

At a point of a plaque, divide the complex tangent space by its null
space.  This gives a complex vector space of dimension $r$ with a positive
Levi form.  The chosen real Lagrangian is complex independent: if its
real span contained both $a$ and $Ja\ne0$, its symplectic form would
vanish on $(a,Ja)$, contrary to Levi positivity.  Its complexified
evaluation, the leaf tangent, and the complexified transverse direction
therefore span the ambient complex tangent space.  This proves the
full-rank assertion.

Let $e_1,\ldots,e_r,f_1,\ldots,f_r$ be a fixed real symplectic frame of
the contact quotient, with the $e_j$ the chosen Lagrangian.  Normalize
the characteristic form by its value on $\partial_s$.  Its pullback is
the fixed contact form, so its exterior derivative on the contact
directions is the fixed symplectic form $\omega$, including as the point
moves along the plaque.  The infinitesimal evaluations of the frame are
holomorphic in the plaque variable.  Use the evaluations of the $e_j$
as a complex frame, and write those of the $f_j$ as the columns of $H$.
The complexified kernel of evaluation consists of $(-H\xi,\xi)$.
It is isotropic for the complexification of $\omega$, since $\omega$
has type $(1,1)$.  The equation
\[
 \omega((-H\xi,\xi),(-H\eta,\eta))
                =\xi^{\mathsf T}(H-H^{\mathsf T})\eta=0
\]
proves symmetry.  Choose the sign of the second symplectic frame so that
$B=\Im H$ is positive definite.  For a vector $a$ in the first real
Lagrangian, its $J$-multiple in real symplectic coordinates is
\[
                  Ja=(-\Re H\,B^{-1}a,B^{-1}a).
\]
Consequently the metric $\omega(a,Ja')$ on that Lagrangian has matrix
$B^{-1}$.

Here is the Taylor calculation, including the mixed terms.  In the
complexified family let $\widetilde\rho(b,z)$ be the pulled-back defining
function, positive on the selected side, and write $b=x+iy$.  On the
real parameter family it is identically zero.  The Legendrian choice
gives $\widetilde\rho_{y_w}(x,z)=0$ identically there, whereas
$a(z)=\widetilde\rho_{y_s}(x_1,z)>0$ after orientation.  Thus every pure
$x_w$ derivative and every $x_wy_w$ derivative occurring in the
weight-two Taylor polynomial vanishes.  Dividing by $a(z)e^2$ gives,
uniformly on compact sets in $(s,w,z)$ near the original plaque,
\[
 \frac{\widetilde\rho(x_1+(e^2s,ew),z)}{a(z)e^2}
       \longrightarrow\Im s-(\Im w)^{\mathsf T}A(z)\Im w.
\]
The positive matrix $A$ is the normalized Levi metric on the chosen
Lagrangian, up to the fixed positive convention in the Levi form and
Taylor coefficient.  Rescale the symplectic frames by that constant;
the preceding computation then gives $A=(\Im H)^{-1}$ exactly.

Finally all infinitesimal evaluations come from a family algebraic in
the projective fibre variable.  Their entries are rational sections on
the plaque variety.  Taking the indicated frame ratios and pulling to
its projective model makes $H$ meromorphic on $L$.  This also proves
the asserted bounded algebraic and analytic parameter dependence.
\end{proof}

\subsection{The first kernel and its whole fibres}

Complexify the family in Lemma~\ref{win:period}, pull back $V$, and take
the open kernel for $b_1=x_1+(e^2s,ew)$.  Denote by $E$ the component
containing $((i,0),z_0)$ for a point $z_0$ of the original regular plaque.
The fixed-weight specialization results show that $E$ is semialgebraic
and locally pseudoconvex, and that it is invariant under real
translations in $(s,w)$ and under $(s,w)\mapsto(a^2s,aw)$, $a>0$.
These actions preserve its component because they are connected actions
starting at the identity.

\begin{lemma}[The whole seed fibre has a Bergman function]\label{win:seed-bergman}
At every $z_0$ in a sufficiently small original plaque patch, the
component of the kernel fibre containing $(i,0)$ is precisely
\[
 T_A=\{(s,w)\in\C^{r+1}:\Im s>(\Im w)^{\mathsf T}A\Im w\},
                    \qquad A=(\Im H(z_0))^{-1}>0.
\]
Its Bergman kernel is strictly positive at every point.  This is a
statement about the entire fibre component, not a relatively compact
part of it.
\end{lemma}

\begin{proof}
Let $K$ be any compact set in the entire scaled $(s,w)$-space.  For
sufficiently small $e$, the points $x_1+(e^2s,ew)$ for $(s,w)\in K$,
and $z$ in a fixed smaller plaque patch, lie in the original regular
defining chart.  The Taylor convergence in the preceding proof holds
uniformly there.  Strict positive points of $q_{z_0}$ have neighborhoods
eventually in the chosen side; strict negative points have neighborhoods
eventually outside it.  A zero of $q_{z_0}$ cannot be in the open kernel,
because every neighborhood contains a strict negative point.  Since
$K$ was arbitrary, this identifies the whole positive component with
$T_A$.  It is connected.  If $r=0$ and both original sides occur, the
two halfplanes are distinct fibre components: the real axis consists
of the original boundary family and is still forbidden.  The component
through $i$ is the upper halfplane in this case as well.

There is an explicit biholomorphism of $T_A$ with a ball.  First set
\[
          S=s+\frac{i}{2}w^{\mathsf T}Aw,
          \qquad W=\frac1{\sqrt2}A^{1/2}w.
\]
An immediate expansion gives
$\Im S-|W|^2=\Im s-(\Im w)^{\mathsf T}A\Im w$.
The map
\[
          (S,W)\longmapsto
          \left(\frac{S-i}{S+i},\frac{2W}{S+i}\right)
\]
then maps the Siegel half-space onto the unit ball in $\C^{r+1}$.
The complex Jacobian $g$ of their composition is nowhere zero, and
the change-of-variables formula yields
\[
                     \int_{T_A}|g|^2
                       =\operatorname{Vol}(\mathbb B^{r+1})<\infty.
\]
Thus every point evaluation on $A^2(T_A)$ is nonzero.  If the fibre
of $E$ has other components, extend $g$ by zero to them; this gives
a holomorphic square-integrable function on that same whole fibre.
No extension from a proper subdomain is being used.
\end{proof}

\begin{lemma}[Exact first epigraph]\label{win:first-epigraph}
There is a connected semialgebraic locally pseudoconvex domain
$V_2\subset L$ such that $H$ is holomorphic on $V_2$, $\Im H>0$ there,
and
\begin{equation}\label{win:epigraph}
 E=\{(s,w,z):z\in V_2,\quad
              \Im s>(\Im w)^{\mathsf T}(\Im H(z))^{-1}\Im w\}.
\end{equation}
In particular $V_2=\{z:((i,0),z)\in E\}$.
\end{lemma}

\begin{proof}
We first classify the fibres away from a proper real algebraic subset
of $L$.  Refine a Boolean algebraic description of $V$ using the
irreducible real polynomial factors of its boundary equations, and
take the factor defining the chosen regular hypersurface patch.
Any other irreducible factor cannot vanish on every generic plaque:
varying all $2r+1$ real transverse parameters would make it vanish on
an open part of the chosen irreducible hypersurface.  The condition
that its restriction to a plaque is identically zero is a finite
algebraic condition on the bounded-degree family.  Exclude its proper
parameter locus.  For the resulting generic plaque, discard the zeros
of the remaining boundary equations, the singular locus of the chosen
equation, and the zeros or poles of the frame determinants and normal
derivative.  These form a proper real algebraic subset of $L$: none
contains the original regular plaque patch.

At any remaining point, the original Boolean description is locally
empty, full, or consists of one or both sides of the chosen equation.
The identities of the Taylor calculation and of $A=(\Im H)^{-1}$
continue as identities of real meromorphic functions on $L$.
Accordingly the full kernel fibre there is empty, full, or one or both
of the two strict quadratic sides.  Each side is connected, being an
epigraph or a hypograph of a quadratic form.  If such a side meets $E$,
it is wholly in its fibre, since a path in that side stays in the same
total kernel component.

Each fibre component is a pseudoconvex tube and hence is convex by
Bochner's tube theorem \cite[Theorem~1.1]{Noguchi2020}.  Moreover a generic component
of $E$ cannot contain a complex affine line.  Indeed, over a small
coordinate polydisk $U\subset L$, the total domain
$E\cap(\C^{r+1}\times U)$ is pseudoconvex in a Stein coordinate
space.  Berndtsson's theorem, with weight zero, says that
\[
                 ((s,w),z)\longmapsto\log K_{E_z}(s,w)
\]
is plurisubharmonic, allowing the value $-\infty$
\cite[Theorem~1.1]{Berndtsson2006}.  The slices here are the actual
unbounded slices; the fibre coordinates have not been truncated.
These local assertions therefore agree on overlapping base boxes.
If the restricted total domain is disconnected, apply the theorem
to its components, or to an exhaustion componentwise.  The Bergman
kernel at a point only uses its fibre component, so this gives the
same function.

A convex domain containing a complex affine line is invariant under
translation in that line direction.  To see this, take the convex
hull of a small ball at any of its points and successively longer
segments in the line; openness then gives each fixed translate of a
smaller ball.  In complex linear coordinates the domain is consequently
a product with $\C$.  Fubini and the fact that an entire square-integrable
function on $\C$ is zero imply that its Bergman space is zero.
An open chamber of such fibres would make the displayed plurisubharmonic
function identically $-\infty$ on a nonempty open subset of connected
$E$.  This contradicts Lemma~\ref{win:seed-bergman}.

At a generic point the quadratic matrix is nonsingular.  Its epigraph
is convex precisely when it is positive semidefinite, and its hypograph
is convex precisely when it is negative semidefinite.  Hence the only
generic nonempty components remaining are the positive definite
epigraph, in $\{\Im s>0\}$, and the negative definite hypograph, in
$\{\Im s<0\}$.  No point of $E$ can have $\Im s=0$: an open
neighborhood of such a point would contain a point over a generic
base with that same imaginary scalar coordinate, contrary to this
classification.  Connectedness and the seed select the positive sign.
Thus $\Im H>0$ on a dense open subset of the projection $V_2$ of $E$.

The meromorphic Cayley matrix
\[
                    C=(H-iI)(H+iI)^{-1}
\]
has operator norm less than one on that dense subset.  It is bounded
where meromorphic and therefore extends holomorphically across its
polar locus.  Its norm is at most one throughout $V_2$.  Equality on
some unit vector would, by the maximum principle applied to a suitable
scalar matrix coefficient, give equality on the connected domain,
contradicting strict contraction on the original plaque patch.  Thus
$\|C\|<1$ everywhere, $I-C$ is invertible, and
$H=i(I+C)(I-C)^{-1}$ is holomorphic with positive imaginary part.
This defines the full epigraph $F$ on the right of
\eqref{win:epigraph}.  The generic classification and openness imply
$E\subset F$, since the non-strict limiting inequalities cannot have
an interior zero, by variation of $\Im s$.

We give the remaining fullness argument, to keep it distinct from the
Bergman argument.  At a generic base point $E_z=F_z$.  Stratify the
exceptional base set.  At a regular codimension-one wall point in
$V_2$, some point of $F_z$ is in $E$ by the definition of the projection.
Every compact path in the connected slice $F_z$ is covered by finitely
many small coordinate boxes compactly contained in $F$.  On a generic
side of the wall all these boxes are available.  The tangency
propagation lemma, Lemma~\ref{holes:tangency-propagation}, applied
successively to these boxes, propagates the one available central
patch along the path.  It uses only transverse parameter disks and
relatively compact fibre patches, so compactness of the whole tube
is not required.  We conclude $E_z=F_z$ at the generic wall points.

The remaining omission $A=F\setminus E$ is relatively closed and
subanalytic of real codimension at least two.  Local pseudoconvexity
and Theorem~\ref{holes:thin-analyticity} make $A$ complex analytic.
It is invariant under all real translations in the fibre coordinates.
For each such direction, analyticity continues this invariance to
small complex translations: restrict its local defining functions to
the complex translation parameter and apply the identity theorem.
Iteration along paths in the connected slice $F_z$ shows that if
$A_z$ is nonempty it is all of $F_z$.  This is impossible for
$z\in V_2=q(E)$.  Hence $A$ is empty.  Finally $V_2$ is the inverse
image of $E$ by $z\mapsto((i,0),z)$, proving its local pseudoconvexity
and the last assertion.
\end{proof}

\subsection{Polynomial normalizations}

The first epigraph reduces the boundary problem to its projective
plaque model. When that model is scaled again, its period matrix
varies with the second scale. The next two lemmas normalize this
matrix while retaining polynomial limits on every bounded box;
this is the compatibility needed for the induction.

We record the finite-dimensional argument used below.  In the phrase
``joint Puiseux expansion,'' a single substitution $t=v^q$ is understood,
after which a fixed power of $v$ times each function is analytic on
a neighborhood of $\{0\}\times U$.  An expansion separately for every
$x\in U$ does not have this meaning.  The generic parameter statement
of Lemma~\ref{prep:puiseux} is used in precisely this joint form.

\begin{lemma}[Finite Taylor extraction]\label{win:finite-taylor}
Suppose $A_t(x)$ and $A_t(x)^{-1}$ have joint Puiseux expansions on
a real parameter neighborhood, with polynomial bounds in $t^{-1}$,
and $R_t(x)$ has the same property and tends to zero locally uniformly.
If $f$ is holomorphic on a fixed complex neighborhood, then any
expression obtained by multiplying $f(x+R_tz)$ on either side by such
Puiseux matrices, and subtracting a similarly bounded constant term,
has a polynomial holomorphic limit on every bounded $z$-box whenever
it is locally bounded on one nonempty complex open $z$-set.  The limit
and convergence are analytic in a generic real center and smooth to
all orders in $z$ on bounded boxes.

In addition, if $R_t$ is invertible and
\[
       C_t(x,u)=R_t(x)^{-1}R_t(x+R_t(x)u)
\]
is bounded on real bounded $u$-boxes, it converges smoothly on those
boxes to a polynomial $G_x(u)$.  After generic restriction,
$\det G_x(u)=1$, and the limiting matrices have polynomial inverses.
\end{lemma}

\begin{proof}
Put $v=t^{1/q}$ with a common denominator for the finitely many
Puiseux expansions.  There are integers $a>0$ and $b\ge0$ such that
$R_t=v^a A(v,x)$, all multiplying factors have norm $O(v^{-b})$,
and their analytic parts extend to a common complex neighborhood.
Taylor-expand $f(x+R_tz)$ through order $m$.  On any fixed bounded
complex $z$-box the remainder after multiplication is
$O(v^{a(m+1)-c})$, where $c$ is a fixed integer bounding the poles
of all factors.  Choose $m$ with $a(m+1)>c$.  The retained part has
a convergent Laurent expansion in $v$, with finitely many negative
terms and coefficients polynomial in $z$.  Boundedness on a nonempty
complex open set forces the coefficient of the lowest negative power
to vanish there and hence identically as a polynomial.  Repeating
removes all negative powers.  The coefficient of $v^0$ is the claimed
polynomial, and the Taylor remainder proves convergence on every
bounded complex box.  Cauchy estimates give derivative convergence.
The same argument with the coefficient parameter $x$ retained gives
the generic local analytic statement.

For the ratio, Taylor-expand the analytic part of
$R_t(x+R_t(x)u)$ in the increment $R_t(x)u$, and then multiply by
$R_t(x)^{-1}$.  Again sufficiently many Taylor terms leave a vanishing
remainder.  Boundedness for real $u$ on an open box kills every negative
Laurent coefficient, because a polynomial vanishing on a real open
box vanishes identically.  This proves even locally uniform
holomorphic convergence on bounded complex $u$-boxes.
Finally, at a generic $x$,
\[
                  \det R_t(x)=v^k\bigl(d(x)+O(v)\bigr),
                         \qquad d(x)\ne0.
\]
Since $R_t(x)u\to0$, the quotient of these determinants at
$x+R_t(x)u$ and $x$ tends uniformly to one.  Thus $\det G_x=1$.
The adjugate formula gives both its polynomial inverse and uniform
boundedness of the inverses on fixed boxes.
\end{proof}

\begin{lemma}[Normalizing a varying Siegel matrix]\label{win:siegel-normalization}
Let $S_2$ be a fixed parameter domain with normalization $R_t^2(x_2)$,
open-kernel component $P_2$, and a fixed interior seed $p_2$.
Let $H_*$ be symmetric and holomorphic on a full neighborhood of $x_2$,
with $\Im H_*>0$ on $S_2$.  Assume the normalizations have joint
Puiseux expansions and polynomial inverse bounds.  Define
\begin{align}
 T_t&=\bigl(\Im H_*(x_2+R_t^2p_2)\bigr)^{1/2},\label{win:T}\\
 H_t(z)&=T_t^{-1}\bigl(H_*(x_2+R_t^2z)
             -\Re H_*(x_2+R_t^2p_2)\bigr)T_t^{-1}.\label{win:Ht}
\end{align}
Then $H_t$ converges on all bounded complex boxes to a symmetric
holomorphic polynomial matrix $\tau$, with $\tau(p_2)=iI$ and
$\Im\tau>0$ on $P_2$.
\end{lemma}

\begin{proof}
The positive square root and its inverse are subanalytic algebraic
operations on positive matrices.  Generic joint Puiseux preparation
therefore applies to $T_t$ and $T_t^{-1}$; positivity of the seed
matrix and the nonzero leading coefficients give their power bounds.
Every compact subset of $P_2$ lies in the normalized $S_2$ for small
$t$, so $H_t$ maps there into Siegel space and $H_t(p_2)=iI$.
Apply the matrix Cayley transform.  Its values lie in the unit matrix
ball and its seed value is zero.  Montel's theorem and the maximum
principle imply bounds for the inverse Cayley transform on compact
subsets of connected $P_2$: a limiting contraction cannot reach the
unit sphere at an interior point while taking zero at $p_2$.
Thus $H_t$ is locally bounded on a nonempty complex open set.

Lemma~\ref{win:finite-taylor} now gives the polynomial limit on all
bounded boxes, not merely within $P_2$.  Its imaginary part is
nonnegative on $P_2$.  For every fixed nonzero complex vector $v$,
the harmonic function $v^*\Im\tau\,v$ is nonnegative and is
$|v|^2$ at $p_2$.  The strong minimum principle makes it positive
everywhere.  This proves strict positivity of the matrix.
\end{proof}

\subsection{The induction and its parameter tests}

\begin{theorem}[Projective windows]\label{win:construction}
Let $Y$ be a connected smooth projective variety and let
$V\subset Y$ be a nonempty connected semialgebraic open subset that
is locally pseudoconvex at its boundary.  The following data exist,
locally at generic centers in totally real parameter boxes.  The
same assertion holds in smooth projective families with additional
real parameters, analytic coefficients, and bounded algebraic
fibre data, when the hypotheses hold fibrewise.  The constructed
data depend analytically on those parameters after generic
restriction and may be complexified locally.

There is a complex box $B\subset\C^d$, a smooth proper projective
family $q:Q\to B$ with connected fibres of dimension $\dim Y-d$,
and a holomorphic evaluation $\mathrm{ev}:Q\to Y$, algebraic with
bounded data in the projective variables and of full rank on a
coordinate patch meeting every fibre.  Its fibre images
$Z_b=\mathrm{ev}(Q_b)$ are irreducible projective varieties of
dimension $\dim Y-d$, and $Q_b\to Z_b$ is proper and surjective.
If $G=\mathrm{ev}^{-1}(V)$, a component $S$ of $q(G)$ and a
relative proper algebraic subset $D\subset Q$ satisfy
\begin{equation}\label{win:container}
               Q_b\setminus D_b\subset G_b\qquad(b\in S).
\end{equation}
The container $D$ is analytic on the full box and proper in every
fibre, after shrinking the box.  The actual holes $Q|_S\setminus G$
are analytic; their images give precisely $Z_b\setminus V$ and
are algebraic in each fibre.  The domain $S$ is locally
pseudoconvex at its parameter boundary.  Its boundary inside $B$
is disjoint from $q(G)$.  All stated shrinkings of $B$ may be
made before the choice of scales and sequences.

There are a fixed interior seed $p$, real block diagonal matrices
$R_t(x)$ tending to zero with polynomial inverse bounds, and a
component $P$ of the open kernel of
\[
                       S_{t,x}=R_t(x)^{-1}(S-x)
\]
at $p$.  A fixed neighborhood of $p$ is eventually in $S_{t,x}$.
The domain $P$ is a finite recursive tower, beginning at a point and
adding a block $(s,w)\in\C\times\C^r$ over a previously constructed
domain $P_2$ by the condition
\begin{equation}\label{win:tower}
      z\in P_2,\qquad
      \Im M(s,w,z)>0,\qquad
      M(s,w,z)=\begin{pmatrix}s&w^{\mathsf T}\\w&\tau(z)\end{pmatrix},
\end{equation}
where $\tau$ is symmetric and polynomial holomorphic and
$\Im\tau>0$ on $P_2$.

For each bounded real $u$-box the ratios
\begin{equation}\label{win:ratio}
           R_t(x)^{-1}R_t(x+R_t(x)u)\longrightarrow G_x(u)
\end{equation}
converge with every derivative.  The matrix $G_x$ is polynomial and
block diagonal, each block depends only on later block coordinates,
$G_x(0)=I$, and every block has determinant one.  Its inverse is
polynomial.  For all sufficiently small real $u$,
\begin{equation}\label{win:orbit}
                           Z\longmapsto u+G_x(u)Z
\end{equation}
is an automorphism of $P$.

Finally the bounded-path localization conclusion holds: any limit
of endpoints of paths in $S_{t,x}$ starting at the seed and staying
in a fixed bounded box belongs to $\overline P$.  Equivalently, on
each fixed bounded box the accessible components have outer
distance to $\overline P$ tending to zero.
\end{theorem}

\begin{proof}
We induct on $\dim Y$.  If the boundary has no ordinary real
hypersurface patch, semialgebraic stratification shows that its
complement has real codimension at least two.  Indeed an open
exterior would be separated from the nonempty open set $V$ by a
stratum of codimension one.  Theorem~\ref{holes:thin-analyticity}
then makes $Y\setminus V$ analytic, and hence algebraic by Chow's
theorem.  Set $d=0$, $Q=Y$, $P=S$ a point, and take this complement
as the container.  In a family, Lemma~\ref{holes:bounded-containers}
selects a containing proper algebraic set with bounded equations;
generic analytic selection of their normalized coefficients gives
a real analytic container.  Complexify those coefficients and
use Lemma~\ref{holes:relative-divisor-container}, shrinking so that
the resulting relative divisor is proper on every fibre of the
full box.  The actual complements need not themselves be selected
as an analytic parameter family.  The same construction is available as a terminal
step as soon as the complement consists only of algebraic holes.
All zero-dimensional assertions have their evident empty-matrix
meaning.

\emph{Composing the projective families.}
Otherwise choose a regular algebraic boundary patch of constant Levi
rank, orient its inward side, and apply
Lemmas~\ref{win:plaques}--\ref{win:first-epigraph}.  This gives the
geometric first kernel $E$ over the smaller projective model $L$,
with projection $V_2\subset L$ and $\dim L=\dim Y-r-1<\dim Y$.
Write $b_1$ for the complex first-family parameters and $x_1$ for
their real center; $e$ will be their scale.  Apply the induction
hypothesis to $(L,V_2)$, keeping $x_1$ and the unused real contact
parameters as external parameters.  Its output is
\[
 q_2:Q_2\to B_2,\qquad D_2,\qquad S_2,\qquad
 (P_2,p_2,R_t^2).
\]
Here $b_2\in B_2$ is the lower-stage parameter, $x_2$ its real
center, and $t$ its scale.  For now $b_2$ is left unscaled.

Replace $D_2$, if necessary, by the full-box proper relative divisor
of Lemma~\ref{holes:relative-divisor-container}.  This only reduces
the available complement and preserves its containment in the open
family.  Complexifying the analytic dependence on $x_1$ composes
the lower family with the first plaque family.  We obtain
\[
 q:Q\longrightarrow B_1\times B_2,
 \qquad \mathrm{ev}:Q\longrightarrow Y,
 \qquad G=\mathrm{ev}^{-1}(V),
\]
where $q$ is smooth, proper, and projective.  The continued divisor,
still denoted $D_2$, is defined on this full box.  All resolutions,
component choices, and box shrinkings here precede the choice of
scales.  We must select one component $S$ of $q(G)$ retaining the
complements of $D_2$, then identify its first kernel.  This will
allow the lower scale $t$ to be introduced without changing the
actual domain being normalized.

On a projective fibre of $Q_2$ over $S_2$, the pullback of $H$ is
constant.  It is meromorphic on that connected projective fibre and
lies in Siegel space off $D_2$.  Its Cayley transform is bounded
there, extends across $D_2$, and is constant on a compact connected
complex manifold.  The resulting symmetric matrix is therefore a
meromorphic function $H_*(b_2)$, with positive imaginary part on
$S_2$; its dependence on the retained first parameters is understood.
Constancy along fibres persists meromorphically on the full connected
box by uniqueness.  Evaluate on a local holomorphic section avoiding
indeterminacy to make $H_*$ holomorphic on a full smaller box about
a generic real center $x_2$.  A proper complex analytic polar set
cannot contain an open subset of the totally real center box, so
this choice is available.

\emph{The fixed parameter domain and its whole fibres.}
We now pass from the limiting fibres of $E$ to the actual fibres of
$G$.  Fix an open chamber compactly inside $S_2$ and let
$b_1=x_1+(i\eta,0)$, $\eta>0$.  Every point of a central projective
fibre off $D_2$ evaluates into $V_2$.  By
\eqref{win:epigraph}, a neighborhood of each such point is eventually
in the inverse image of $V$.  Properness gives the following useful
closed-set formulation: every accumulation point, as $\eta\downarrow0$,
of a hole on a compact parameter subbox belongs to $D_2$.  Otherwise
a convergent sequence of holes would end in one of the just-described
eventually included neighborhoods.

Here the limiting statement holds for all approaches through the
upper half-disk, not just for the vertical ray.  Write its scalar
coordinate as $u+i\eta$, recenter the real first parameter at $u$,
and put $e=\sqrt\eta$.  In normalized coordinates the point is
again $(i,0)$.  On a compact fibre patch outside $D_2$, this point
has a positive distance from the limiting forbidden set.  Apply
Lemma~\ref{prep:generic-distance-convergence} to the capped distances
of the normalized forbidden sets, including the remaining real
parameters and the compact fibre variables jointly.  Use one fixed
bounded seed box and finitely many compact projective coordinate
charts.  The finite-test part of that lemma supplies one open real
center neighborhood on which convergence is uniform on compact
parameter and test boxes.  Now choose an accuracy smaller than half
the positive separation of the patch under consideration.  The
same separation holds for all sufficiently small $e$ and nearby
real centers.  The threshold may depend on the patch; no positive
lower bound for the separation near $D_2$ is required.  Every sequence of upper
half-disk holes approaching a point off $D_2$ would violate this
separation on a compact patch about its limit.  Thus all such
accumulation points belong to $D_2$, as required by the boundary
trapping hypothesis.  The container has a holomorphic continuation
on the full box.
Lemma~\ref{holes:boundary-trapping}, applied to the scalar upper
half-disk, therefore traps the holes in that continuation on the
indicated side.  For the remaining first parameters use
Lemma~\ref{holes:interior-trapping}: their real boxes are totally
real and nonpluripolar.  This gives actual containment on a nonempty
complex open parameter region.  Choose the component $S$ of $q(G)$
reached there.  Proposition~\ref{holes:projective-propagation} makes
the actual holes analytic throughout $S$, and
Corollary~\ref{holes:container-propagation} keeps them inside this
same full-box divisor $D_2$.
Their proper images are precisely $Z_b\setminus V$;
these are compact analytic and hence algebraic.  The container is
proper fibrewise after shrinking at a generic center.
Corollary~\ref{holes:section-test} now gives local pseudoconvexity of $S$
at its boundary.  Since $q(G)$ is open, a boundary point of one of
its components cannot belong to any of its components; this proves
the asserted exclusion from $q(G)$.

For the first-kernel comparison below, make one further full-box
restriction.  The total evaluation $Q_2\to L$ is generically a
submersion.  Choose at the generic center a point off $D_2$ where
its differential has full rank.  Smoothness of the projective
family gives a holomorphic section $\sigma$ through this point.
Shrink the box so that $\sigma$ is everywhere off the continued
container and the evaluation retains full rank on a fibre-coordinate
patch around $\sigma$.  Include the first parameters in this
holomorphic choice.  The differential includes both base and
projective fibre variables; it need not have full rank in the
fibre directions alone.

After these box restrictions, retain the selected connected
components $S$ and $S_2$.  Fix $b_2^0$ in the retained chamber of
$S_2$.  For a compact path in $S_2$ starting at $b_2^0$, choose
local sections of $Q_2$ off
$D_2$ along a finite covering of the path.  Their evaluations lie
in $V_2$.  The strict inequalities in \eqref{win:epigraph} and
neighborhood eventual inclusion put their small first-scale
displacements in $q(G)$, with overlaps joining them to the region
used to choose $S$.  Thus every such displaced path lies in this
same $S$ for all sufficiently small first scales.  The threshold
may depend on the compact path, but $S$ is fixed before either
scale tends to zero.

\emph{The first kernel in parameter space.}
Let $\mathcal K_1$ be the full first parameter kernel of this fixed
$S$, with $b_1=x_1+(e^2s,ew)$ and $b_2$ unscaled.  Its component
$A$ through $(i,0,b_2^0)$ is exactly
\begin{equation}\label{win:first-parameter-kernel}
 A=\left\{(s,w,b_2):b_2\in S_2,\quad
       \Im\begin{pmatrix}s&w^{\mathsf T}\\w&H_*(b_2)\end{pmatrix}>0\right\}.
\end{equation}
The compact-path argument and strict neighborhood inclusion in
\eqref{win:epigraph} put the entire displayed domain in this
component.  For the reverse inclusion, a parameter neighborhood in
$\mathcal K_1$ must give a neighborhood in the geometric kernel $E$.
Use the section $\sigma$ and evaluation patch prepared above to
make this an open evaluation test.

Suppose a neighborhood belongs eventually to the first normalized
parameter domain.  By \eqref{win:container} every point in this
evaluation patch is then in the inverse image of $V$.  The
submersion theorem, on a smaller fixed patch, sends these
neighborhoods to neighborhoods in the first plaque family.
The openness here is uniform after normalization.  In local
coordinates the normalized evaluation has the form
\[
 (s,w,b_2,v)\longmapsto
 (s,w,F(x_1+(e^2s,ew),b_2,v)),
\]
where $v$ is a projective fibre coordinate.  The first coordinates
are unchanged, and $D_{(b_2,v)}F$ has a uniformly bounded right
inverse on a smaller patch around $\sigma$.  The remaining
derivative block tends to zero because it contains the factors
$e^2,e$.  The block triangular differential consequently has a
uniformly bounded right inverse on each compact normalized path
neighborhood.  The quantitative local submersion theorem gives
the needed neighborhoods of a fixed positive size.  In the limit
this is a neighborhood test for membership in the full first
geometric kernel.  Along a path $(s(a),w(a),b_2(a))$ from
$(i,0,b_2^0)$ in the selected component,
the tested points
\[
       (s(a),w(a),\mathrm{ev}_2(\sigma(x_1,b_2(a))))
\]
form one continuous path in that kernel.  At the starting point,
the section is off $D_2$ over $S_2$, so its evaluation is in $V_2$
and the tested point is in $E$.  The entire tested path therefore
stays in $E$.  There are no changes of section in this argument.
Its base values stay in $q_2(\mathrm{ev}_2^{-1}(V_2))$ and, by
connection to the seed, in its component $S_2$.  Equation
\eqref{win:epigraph} gives precisely the strict matrix test in
\eqref{win:first-parameter-kernel}.  This proves the reverse
inclusion.  The same rank argument, followed by the first family's
full-rank evaluation, proves full rank of the final evaluation on
a patch meeting every fibre in a smaller box.  Its differential is
then injective also on the fibre tangent.  The irreducible image
therefore has dimension $\dim Q_b=\dim Y-d$; properness of $Q_b$
gives the remaining fibre-image assertions.

\emph{The second specialization and its marked component.}
We have obtained a fixed domain $S$ and identified its first marked
kernel $A$.  Now normalize $H_*$ by \eqref{win:T}--\eqref{win:Ht}.
In the first block use $b_1=x_1+(e^2s,eT_tw)$; in later blocks use
$b_2=x_2+R_t^2z$.  Congruence of the matrix inequality by
$\operatorname{diag}(1,T_t^{-1})$ identifies its first-stage test
with
\[
              \Im\begin{pmatrix}s&w^{\mathsf T}\\w&H_t(z)\end{pmatrix}>0.
\]
Lemma~\ref{win:siegel-normalization} gives its limit
\eqref{win:tower}.  Apply the above $t$-dependent affine normalization
to $\mathcal K_1$ and $A$ to obtain $\mathcal K_{1,t}$ and $A_t$.
The seed $(i,0,p_2)$ lies in $A_t$ for every small $t$, because
$H_t(p_2)=iI$.  Consider a compact path from this seed in the
marked component of the open kernel of $\mathcal K_{1,t}$.
A compact connected tube of small neighborhoods around this path
is eventually contained in $\mathcal K_{1,t}$.  Since that tube
contains the seed, it is contained in its component $A_t$ for all
small $t$.  Its projection supplies neighborhoods eventually in
the normalized $S_2$.  The projected path consequently belongs
to the second open kernel and, by connection to $p_2$, to $P_2$.
The limit matrix is nonnegative on the tube.  With its lower block
positive on $P_2$, varying $\Im s$ shows that a zero Schur
complement cannot be an interior point.  Thus the limiting matrix
test is strict.  Conversely, compact neighborhoods in the strict
tower inequalities are eventually in $A_t$, by kernel inclusion
for $S_2$ and locally uniform convergence of $H_t$; their paths
connect them to the seed.  This proves the equality of marked
iterated kernels.  In particular, other components of the full
first kernel cannot enter this marked component when the second
scale tends to zero.  Only compact-path neighborhood inclusion
is used here, not the later accessible-path localization theorem.

\emph{A single scale.}
Apply the closed-set power diagonalization theorem,
Lemma~\ref{prep:power-diagonalization}, to these fixed definable
forbidden-set tests and their complements.  It supplies an integer
$N$ such that replacing $e$ by $t^N$ preserves the iterated open
kernel on every fixed bounded box, with the seed component just
identified.  Accuracy and box radius are fixed parameters in that
theorem; the exponent is uniform in those parameters whereas the
threshold in $t$ is allowed to depend on them.  Increasing $N$
preserves this conclusion.  We shall do so finitely many times
below.  The final normalization is
\begin{equation}\label{win:Rt}
          R_t=\operatorname{diag}
                     (t^{2N},t^NT_t,R_t^2).
\end{equation}
It tends to zero and has polynomial inverse bounds.  A compact
ball about $(i,0,p_2)$ supplies the fixed seed neighborhood.

\emph{Recentering ratios and localization.}
We verify the ratio condition, including the square-root order.
By induction the later-block ratio $C_t^2(u_2)$ is bounded and
converges polynomially.  Replacing the later center by
$x_2'=x_2+R_t^2u_2$ makes its seed evaluation point equal to
\[
         x_2+R_t^2\bigl(u_2+C_t^2(u_2)p_2\bigr).
\]
This is a bounded normalized point for bounded $u_2$.
Polynomial bounded-box convergence of $H_t$ therefore bounds
\[
        T_t^{-1}(T_t')^2T_t^{-1}
          =(T_t^{-1}T_t')(T_t^{-1}T_t')^{\mathsf T}.
\]
It follows that $T_t^{-1}T_t'$ is bounded.  No commutativity of the
positive square roots is asserted or needed.

All functions involved have joint analytic Puiseux expansions near
generic centers.  Their derivatives of any fixed finite order
have power bounds.  The shift of $x_1$ is
$O(t^N\|T_t\|+t^{2N})$ on bounded boxes.  Choose $N$ larger than
the finitely many pole orders in the lower-scale data, their
inverses, and the Taylor remainders used here.  Its effect on those
data and on the square-root ratio then tends to zero.  The
normal-coordinate ratio is identically one, and the tangential
ratio has a limit depending only on $u_2$.  Apply
Lemma~\ref{win:finite-taylor} to obtain polynomial convergence with
all derivatives.  Applying its determinant argument separately
to $T_t$ and to each old block shows that every limiting block has
determinant one.  The later-variable dependence in older blocks
is the induction hypothesis.  This proves \eqref{win:ratio} and
all its asserted moderation properties.

Finally use Lemma~\ref{prep:generic-distance-convergence} for the
closed forbidden-set tests in the normalized coordinates.  At a
generic real center the distance limits on each fixed box are
unchanged by all centers tending to that center.  A separate
neighborhood and threshold are permitted for each fixed accuracy.
The exact change of center in the original fixed set $S$ is
\[
 x+R_t(x)\bigl(u+C_t(x,u)Z\bigr)
       =x+R_t(x)u+R_t(x+R_t(x)u)Z.
\]
Taking neighborhood limit tests on both sides and using the
bounded inverse ratios proves invariance of the full kernel
under $Z\mapsto u+G_x(u)Z$.  For small $u$, the image of the seed
stays in the same component; hence this is an automorphism of
$P$.  The bounded-path conclusion is now Theorem~\ref{loc:theorem},
applied to the fixed pseudoconvex domain $S$, the seed
ball, and the ratios just proved.  This completes the induction.
Here $D_t^x=S_{t,x}$ in the notation of the localization theorem.
\end{proof}

\subsection{Geometry of the limiting tower}

\begin{proposition}\label{win:tower-properties}
The domain $P$ of Theorem~\ref{win:construction} is diffeomorphic
to $\R^{2d}$, where $d=\dim_{\C}P$.  In particular it is
contractible.  It is Stein and biholomorphic to a bounded domain
of the same complex dimension.  Complete real holomorphic vector fields from
the identity component of its automorphism group span every
complex tangent space over $\C$.
\end{proposition}

\begin{proof}
At one stage of the tower put $B(z)=\Im\tau(z)$.  The Schur
complement rewrites its condition as
\[
       \Im s>(\Im w)^{\mathsf T}B(z)^{-1}\Im w.
\]
The positive slack
\[
 \lambda=\Im s-(\Im w)^{\mathsf T}B(z)^{-1}\Im w
\]
gives smooth coordinates $(z,w,\xi,\eta)$ on the whole stage,
where $\xi=\Re s$ and $\eta=\log\lambda$.  Their inverse recovers
$s=\xi+i\bigl(e^\eta+
(\Im w)^{\mathsf T}B(z)^{-1}\Im w\bigr)$.
Thus this stage is diffeomorphic to
$P_2\times\C^r\times\R^2$.  Induction, starting from a point,
proves the claimed diffeomorphism and contractibility.

For Steinness, let $M_j$ be the polynomial matrices in all blocks.
The function
\[
             \Phi(Z)=|Z|^2-\sum_j\log\det\Im M_j(Z)
\]
is strictly plurisubharmonic on $P$.  The matrix entries have
harmonic imaginary parts, and differentiating $-\log\det$ gives
a nonnegative Levi form, equivalently the squared matrix norm
of the derivative conjugated by the positive inverse square
root.  At a finite boundary point at least one positive matrix
becomes singular, so $\Phi\to+\infty$.  At infinity its negative
terms are at worst logarithmic in $|Z|$, because the $M_j$ are
polynomial, whereas $|Z|^2$ dominates them.  Thus $\Phi$ is an
exhaustion and $P$ is Stein.

For a bounded realization use one block at a time.  Since
$\Im M>0$, the matrix $K=I-iM$ has Hermitian part greater than
$I$, and hence $\|K^{-1}\|\le1$.  Write its first column as
$(a,b)^{\mathsf T}\in\C\times\C^r$.  The lower-right block
$D=I-i\tau(z)$ is invertible, and the Schur complement formula
gives $a\ne0$.  The equations for that column recover the original
coordinates holomorphically:
\[
       w=\frac{Db}{ia},\qquad
       s=\frac{i(1-a+iw^{\mathsf T}b)}{a}.
\]
Thus $(s,w,z)\mapsto(a,b,\psi_2(z))$, with an inductively obtained
bounded realization $\psi_2$, is injective, bounded, and has a
holomorphic inverse on its image.  It uses exactly $r+1$ new
coordinates and its image is open.  This proves the
same-dimensional bounded realization.

The automorphism group of a bounded domain is a real Lie group
\cite[Proposition~3.1 and Theorem~6.2]{Kobayashi1967}.
Differentiate the curves \eqref{win:orbit} at $u=0$ to obtain
complete real holomorphic vector fields.  In the block order of
the theorem their values have identity diagonal: differentiation
in a coordinate of block $j$ contributes its coordinate unit
vector in that block, nothing in later blocks, and possibly
vectors in earlier blocks, because $G_k$ only depends on blocks
later than $k$.  This triangular matrix has determinant one at
every point.  The resulting fields therefore span the complex
tangent space everywhere.
\end{proof}
